\newpage
\section{Drone Detection Using Sensors}
\label{sec:detection}
Detection must \emph{find} a small object in the sky, \emph{classify} it as
a drone or not, for example a bird, and give the \emph{direction} to it. A
target can be found when it covers a few pixels, but it can only be
classified reliably when it covers tens of pixels. Detection is therefore
split into a cheap first step that flags candidates and a classifier that
runs only on them.

\subsection{Sensors}
Each defender has a camera on a gimbal that scans its outward sector
(Section~\ref{sec:base_surveillance}). The baseline uses the simulated
platform's camera (Table~\ref{tab:assumptions}); the target system's camera
is chosen in Section~\ref{sec:camera-choice}.
Classification needs more pixels than detection. Johnson's criteria give
about eight pixels for recognising an object \cite{johnson1958analysis},
and a detector such as YOLO needs more in practice, so classification is
assumed to need about 20~pixels, a value that will be measured in the
simulation. By
Equation~\eqref{eq:rdet}, a camera then classifies a target at a fifth of
its detection range: about 10~m with the platform's camera and about
136~m with a 20$^\circ$ camera at 1920~px.

A thermal camera on the gimbal is an optional extension that could be tried
later: a drone's warm motors stand out against a cold sky, also at night
and against trees, where an RGB camera struggles. LiDAR is not used, since
a small drone at long range falls between the beams. A real system would
also have a radar for early warning, but Gazebo has no radar sensor.

\subsection{Baseline: Dark Object Against the Sky}
The baseline groups pixels that differ clearly from the surrounding sky
into connected regions, and reports regions of a plausible size as
candidates. It needs no training, but it cannot tell a drone from a bird
and fails when the background is not sky. Every candidate seen in $k$
consecutive images is treated as a drone.

\subsection{Improvement: YOLO26}
To detect and classify hostile drones, we will use YOLO26 from Ultralytics,
the newest version of the YOLO model family
\cite{jocher2026ultralyticsyolo26unifiedrealtime}. YOLO26 looks at an image
once and directly outputs where objects are (bounding boxes), what they are
(classes), and how confident it is, which makes it fast enough to run in
real time. Tiny drones are hard for standard YOLO models
\cite{zhai2023yolodrone}. Of YOLO26's five sizes, from nano (n) to extra
large (x), we use the two smallest, since computing power on board is
limited \cite{wang2026edge}. YOLO26n runs on
the drone, on enlarged crops of the image around the candidates, which lets
it see small targets \cite{akyon2022sahi}. YOLO26s is a candidate for the
stronger classifier at the base station (Section~\ref{sec:base-station}).

Detections are linked between images with
ByteTrack \cite{zhang2022bytetrack}, so a target keeps its identity when it
is missed in some images. A drone starts flying towards a target as soon as
YOLO26n classifies it as a drone with a confidence above $\tau$ in $k$
consecutive images (Section~\ref{sec:intercept}).

\subsection{Training Data and Fine-Tuning}
\textbf{Data collection.} Most of the training data is generated in Gazebo,
where every image can be labelled automatically. A script places drone and
bird models at random positions around a defender and points its gimbal
camera at random directions, varying the distance, altitude (20--120~m),
background (clear sky, clouds, trees and ground) and lighting. The distances
are chosen so that the targets span from a few pixels to about 100 pixels
across, with most images near the 4- and 20-pixel limits that matter for
detection and classification. Each object's bounding box is computed from
its true position and the camera's pose, or read directly from Gazebo's
bounding-box camera sensor, which also handles objects hidden behind trees.
Images with no drone or bird are included as well, so that the model learns
to ignore clouds and trees. The images are split into training, validation
and test sets by simulation run rather than by image, since consecutive
images from the same run are nearly identical and would otherwise make the
test results look better than they are.

\textbf{Fine-tuning.} Training starts from Ultralytics' YOLO26n and YOLO26s
weights pretrained on COCO, whose backbones already extract general visual
features. The classification part of the detection head, which for COCO has
one output per each of 80 classes, is replaced with one that has only two
outputs, \emph{drone} and \emph{bird}, and is trained from scratch, while the
rest of the network keeps its pretrained weights. Everything else in an image,
such as clouds, trees and buildings, is treated as background, so the model
learns to ignore it from the training images. Training is done in two steps.
First, the models are fine-tuned on the Seraphim drone dataset
\cite{seraphim}, about 83\,000 real drone images, so that they learn what
drones look like in many settings. Second, they are fine-tuned on the Gazebo
images, with part of the Seraphim images mixed in so that the models do not
forget them. The Drone-vs-Bird dataset \cite{coluccia2021drone}, real
videos of small drones with birds as distractors, could be added if birds
turn out to be hard to reject. Since YOLO26n runs on crops of the camera
image, it is trained on crops of the same size, so that a target has the
same size in pixels during training as during use \cite{akyon2022sahi}.
Training runs on a GPU with the standard Ultralytics augmentations, such as
changes in brightness and scale, and stops when the result on the validation
set no longer improves. The confidence threshold $\tau$ is then chosen on
the validation set so that almost no drones are missed. All results are
reported on a test set of held-out Gazebo runs, since the system has to
perform in the simulation.

\subsection{Direction and Identification}
The centre of a bounding box gives the direction to the target, from the
camera's intrinsic parameters and the gimbal's orientation
\cite{hartley2004multiple}. One camera does not measure the range, but the
width of the box gives a rough estimate: assuming the target has the size
$s$ in Table~\ref{tab:assumptions}, a target that covers $p$ pixels is at
range $r \approx f s / p$, as in Equation~\eqref{eq:rdet}. This estimate is
uncertain, since the hostile drone's true size is unknown and an error of
one pixel is a large part of a target only a few pixels wide. The drone
therefore sends the base station an estimated position together with its
uncertainty, which is small across the line of sight and large along it.
The base station's Kalman filter combines these estimates with those of
other drones (Section~\ref{sec:base-station}). Once two or more drones see
the target, triangulation (Equation~\eqref{eq:triangulation}) gives an
accurate position; until then, the single-drone estimate fills the gap.
The base station also knows every defender's position from its telemetry,
so a detection within a few metres of a defender's reported position is
classified as friendly. Only the remaining detections can be hostile.

\subsection{Evaluation}
The dark-object baseline and YOLO26 are compared on the same seeded
scenarios, in which a hostile drone and birds approach the base from random
directions. A detection counts as correct if its centre lies within a few
pixels of the target's true position projected into the image, since box
overlap is unstable for targets only a few pixels wide. The measures are:
\begin{itemize}
    \item \textbf{Missed drones:} the share of hostile drones that reach the
          ring undetected, the most important measure.
    \item \textbf{Detection probability versus range:} the share of hostile
          drones detected in each range interval, which shows where detection
          starts to fail.
    \item \textbf{Birds flagged as drones:} the share of birds seen by a
          camera that are classified as drones.
    \item \textbf{Background false alarms:} detections of drones in images
          with no drone or bird, per image.
    \item \textbf{Recall versus target size:} the share of drones detected in
          held-out Gazebo images, grouped by their size in pixels, which tests
          the assumptions of 4 pixels for detection and 20 pixels for
          classification.
    \item \textbf{Inference time} of YOLO26n and YOLO26s per image on a CPU.
\end{itemize}


\subsection{Interface}
\label{sec:detection-interface}
The detection component produces one \emph{detection report} for each
target it follows (Table~\ref{tab:detection-report}). The report is used
in three places: on the drone, every image, by the gimbal controller and
the guidance law (Section~\ref{sec:intercept}); by the scan, which stops
when a target is found (Section~\ref{sec:base_surveillance}); and at the
base station, every $\Delta t$ (Section~\ref{sec:target-assumptions}),
together with an image crop of the target (Section~\ref{sec:base-station}).

\begin{table}[H]
\centering
\small
\begin{tabular}{|l|p{0.62\textwidth}|}
\hline
\textbf{Field} & \textbf{Content} \\
\hline
Time & When the image was taken \\
Drone & Which drone saw the target \\
Track & Identity of the target over time, from ByteTrack \\
Stage & \emph{Candidate} (first step) or \emph{classified} (YOLO26) \\
Class and confidence & \emph{Drone} or \emph{bird}, with the model's confidence \\
Direction & Unit vector from the drone to the target, in the world frame,
            with its uncertainty \\
Range & Rough range from the target's size in the image, with its
        uncertainty \\
Bounding box & Position and size in pixels; used only on the drone, to
               centre the target in the image \\
\hline
\end{tabular}
\caption{Contents of a detection report (draft).}
\label{tab:detection-report}
\end{table}

The direction is given in the world frame rather than in pixels, so that
the components that use it do not depend on which camera is chosen
(Section~\ref{sec:camera-choice}). Direction and range are kept apart
rather than combined into a single position, since the direction is
accurate while the range is not, and the guidance laws need only the
direction. Whether a detection is a friendly drone is decided at the base
station, which knows every defender's position, and is therefore not part
of the report.
% TODO: decide message format (e.g. a custom ROS 2 message or
% vision_msgs/Detection2DArray), the rate on the drone, and how the
% uncertainties are represented.